\section{讨论与结论}

\subsection{收益兑现与风险配置}

完整历史比较呈现出收益水平与实现路径的差异。V3.2 与原始系统的累计收益接近，回撤及月度下行偏差较低，期末未退出头寸减少。风险预算和退出规则为交易设置了损失约束，并在趋势延续时调整盈利保护，但也可能缩短持有有利趋势的时间。熊市窗口中收益与回撤同时下降，说明这种变化存在机会成本。

确认加仓使初始突破承担较小资金风险，并在第二次同向突破后扩大仓位。确认与未确认交易的利润分层符合这一资金配置机制，也解释了负等交易均值与正账户利润的共存。由于利润集中于确认交易，漏加仓、部分成交和动态止损偏离会直接影响策略收益来源。

版本比较还显示风险约束存在不同取舍。V3.1 的最大百分比回撤和赢家集中度低于 V3.2，CAGR/回撤比更高；V3.2 的累计收益及 Profit Factor 更高。两者分别代表较低历史回撤与较高趋势收益的配置，不能仅按版本先后判定优劣。

\subsection{研究局限与后续验证}

本研究属于开发样本上的完整系统比较。第二节所述模型选择偏差限制了结果的外推；原始系统与 V3.2 的杠杆、仓位及退出规则同时存在差异，亦无法据此单独识别每条规则的因果贡献。分窗交易数较少、重采样结果不稳定及收益集中于少数趋势交易，进一步增加了估计不确定性。

执行方面，月度波动指标基于已实现权益，不含持仓期间的全部浮动盈亏。K 线模拟不能完整反映滑点、盘口、资金费率、部分成交和清算路径。BTC 为主要研究资产，其他币种迁移未获得稳定支持，因此结论限于本文样本和执行假设。

后续验证将使用固定规则的 Bybit Demo 前向样本，记录收益、回撤、确认成交率、订单偏离及费用，并补充最终版本成本压力测试和滚动窗口分析。V3.1 可作为共同起点的比较版本。前向结果不纳入本文现有历史统计。

\subsection{结论}

本文考察规则化退出与确认加仓对 BTC 永续波动突破策略收益兑现的影响。在所用完整历史样本中，V3.2 的累计收益与原始系统接近，最大百分比回撤相对降低约 28.6\%，月度下行偏差相对降低约 48.1\%，总体月度收益年化波动率仅小幅下降。其主要改善体现为下行风险压缩和正常规则退出，而不是累计权益持续领先。

确认加仓改变了试探与确认状态的资金权重，历史利润主要由确认交易贡献，同时伴随赢家集中和执行依赖。综合结果支持将退出设计与资金分配作为波动突破策略评价的核心维度，其经济效果仍需结合成本与前向样本进一步检验。
